\documentclass[12pt]{article}
\usepackage[letterpaper,margin=0.75in]{geometry}
\usepackage[T1]{fontenc}
\usepackage{lmodern}
\usepackage{microtype}
\usepackage{booktabs}
\usepackage{array}
\usepackage{enumitem}
\usepackage{amssymb}
\usepackage{hyperref}
\usepackage{setspace}
\setstretch{1.0}
\setlength{\parindent}{1em}
\setlength{\parskip}{0pt}
\setlength{\tabcolsep}{5pt}
\setlist{nosep,leftmargin=1.4em}
\hypersetup{
  hidelinks,
  pdftitle={Newlines are scope: Negation lost in clinical-trial text normalization},
  pdfauthor={},
  pdfsubject={Anonymous PLC 51 abstract},
  pdfkeywords={}
}
\pagestyle{plain}
\renewcommand{\refname}{References}
\newcommand{\ex}[1]{\textit{#1}}

\begin{document}

\begin{center}
{\bfseries Newlines are scope: Negation lost in clinical-trial text normalization}
\end{center}
\vspace{-0.65em}

Eligibility criteria determine who may enter a clinical trial, but they are
distributed as free text that must be converted into computable constraints.
Negation is consequently central to clinical information extraction, from
NegEx (Chapman et al. 2001) to Criteria2Query (Yuan et al. 2019; Fang et al.
2022) and recent trait-decomposition pipelines (Chen et al. 2026). These
approaches model lexical cues, entities, and logical relations. We identify a
prior representational problem: a criterion's line boundary can be its only
overt scope boundary. ClinicalTrials.gov preserves Markdown-like bullets, but
newline-flattening turns sibling criteria into one punctuation-delimited
string. We ask how often this lossy normalization creates a possible negation
span that contradicts the list structure.

\textbf{Data and method.}
We queried the public ClinicalTrials.gov API on November 5, 2026 for the 250
most recently updated recruiting trials in each of four domains: cancer, heart
disease, diabetes, and depression. After cross-query deduplication, the sample
contained 989 trials; 970 had explicitly bulleted or numbered criteria. We
identified 17,673 list units, joining unmarked wrapped lines to their preceding
bullet before analysis. This avoids treating display wrapping as structure.

We located a conservative set of clinical negation cues: \ex{no, not,
without, never, neither, cannot, can't, must not, may not, absence of, free of},
and \ex{negative for}. A \textsc{boundary-sensitive} token is a cue in a list
unit that (i) lacks final punctuation and (ii) is followed by another unit in
the same inclusion or exclusion section. A line-aware representation closes
scope at the list boundary. A punctuation-only baseline deletes the boundary
and holds scope open until \ex{., !, ?, ;}, or \ex{:}. This diagnostic measures
a structural failure opportunity, not an observed error by a particular
system.

Example (1), from NCT07318805, shows three sibling exclusion criteria. The
direct object of \ex{controlled} ends in (1a); neither (1b) nor (1c) is in the
scope of \ex{cannot}. After newline deletion, however, no punctuation separates
the three predicates.

\begin{enumerate}[label=(\arabic*),leftmargin=2.2em,itemsep=0pt,topsep=2pt]
\item \begin{tabular}[t]{@{}ll@{}}
a. & \ex{Hypertension that cannot be controlled by optimal medical therapy}\\
b. & \ex{Renal impairment}\\
c. & \ex{Hepatic dysfunction}
\end{tabular}
\end{enumerate}
\vspace{-0.25em}

\textbf{Results.}
Of 17,673 criteria, 2,084 (11.8\%) contained a cue. Among cue-bearing units,
892 (42.8\%) were boundary-sensitive (Table~\ref{tab:domains}). With newlines
removed, the punctuation baseline crossed a median of four subsequent criteria
and 67 words before reaching a terminator. Rates differed by clinical domain,
$\chi^2(3)=28.05$, $p<.001$, Cramer's $V=.116$: oncology had the highest rate
(47.7\%) and psychiatry the lowest (32.5\%). Oncology's longer spill spans
(median five criteria and 120 words) show that the effect is not merely one
adjacent short phrase. In a deterministic manual audit of 40 events (10 per
domain), 36 (90\%) were independent sibling constraints; four were coordinated
or nested continuations. This motivates retaining indentation and overt
continuation markers rather than treating every newline identically.

\begin{table}[h]
\centering
\caption{Boundary-sensitive negation cues by domain. ``Trials'' includes only
records with structured list units.}
\label{tab:domains}
\vspace{-0.35em}
\begin{tabular}{@{}lrrrr@{}}
\toprule
Domain & Trials & Units & Cue units & Boundary-sensitive\\
\midrule
Oncology   & 241 & 5,417 & 975 & 465 (47.7\%)\\
Cardiology & 245 & 4,057 & 333 & 140 (42.0\%)\\
Diabetes   & 240 & 3,957 & 354 & 150 (42.4\%)\\
Psychiatry & 244 & 4,242 & 422 & 137 (32.5\%)\\
\midrule
Total      & 970 & 17,673 & 2,084 & 892 (42.8\%)\\
\bottomrule
\end{tabular}
\end{table}
\vspace{-0.6em}

\newpage
\textbf{Analysis.}
In this register, typography supplies a constituent boundary. Eligibility
criteria are telegraphic predicates whose subject (the potential participant)
and modal force (must/must not satisfy) are inherited from the section heading.
A same-level bullet packages one constraint under that shared frame. The line
break is therefore not decorative whitespace: it closes the current constraint
before the next predicate is evaluated. Flattening (1) erases the evidence that
\ex{renal impairment} is a separate exclusion, not an argument or continuation
of \ex{cannot be controlled}.

We propose \textsc{boundary-aware scope preservation}. Before lexical or neural
analysis, each same-level criterion boundary is converted to a typed sentinel,
\texttt{<CRIT>}; ordinary wrapped lines remain inside the unit, while indented
sub-bullets attach to their parent. Negation scope may not cross \texttt{<CRIT>}
unless an overt continuation marker (e.g., \ex{and} introducing a coordinated
subcriterion) licenses attachment. This preprocessing change eliminates all
892 mechanical spill opportunities by construction while leaving within-unit
negation to existing systems. It is compatible with rule-based scope systems,
semantic parsers, and language-model tokenizers.

The finding complements annotated resources such as the 1,006-document Leaf
Clinical Trials Corpus (Dobbins et al. 2022). Those resources richly represent
entities and relations, but the input serialization itself remains part of the
semantic problem. Recent work decomposes complex criteria into atomic traits
and explicitly annotates negation (Chen et al. 2026); our result identifies a
cheap, upstream boundary signal for that decomposition. It also offers a
controlled benchmark: pair each original list with a newline-flattened version
and test whether a parser preserves sibling polarity and logical independence.

Three predictions follow. First, boundary preservation should improve
precision most in oncology, where negation cues and unpunctuated lists are most
dense. Second, models trained on flattened corpora should be less robust to
original list formatting than models trained with typed boundaries. Third,
human judgments should assign narrower scope to the formatted version of
minimal pairs like (1), even when lexical content is held constant. The current
cue-based study is not a gold-standard error analysis; rather, it quantifies
how frequently a commonplace normalization destroys a linguistically licensed
scope boundary. The manual audit was small and single-reviewed; clinician
annotation remains necessary. In a safety-relevant medical register, preserving
the original boundary should nevertheless be the default.

\vspace{-0.3em}
\begin{thebibliography}{9}
\setlength{\itemsep}{0pt}
\setlength{\parskip}{0pt}

\bibitem[Chapman et al. 2001]{chapman2001}
Chapman, Wendy W., Will Bridewell, Paul Hanbury, Gregory F. Cooper \,\& Bruce G. Buchanan. 2001.
\newblock A simple algorithm for identifying negated findings and diseases in discharge summaries.
\newblock \textit{Journal of Biomedical Informatics} 34(5). 301--310.

\bibitem[Chen et al. 2026]{chen2026}
Chen, Hongyu, Lingfei Qian, Xing He et al. 2026.
\newblock Clinical trial eligibility criteria decomposition and parsing with large language models.
\newblock \textit{Studies in Health Technology and Informatics}. doi:10.3233/SHTI250928.

\bibitem[Dobbins et al. 2022]{dobbins2022}
Dobbins, Nathan J., Taylor Mullen, Ozgur Unal \,\& Meliha Yetisgen. 2022.
\newblock The Leaf Clinical Trials Corpus: A new resource for query generation from clinical trial eligibility criteria.
\newblock \textit{Scientific Data} 9. 490.

\bibitem[Fang et al. 2022]{fang2022}
Fang, Yilu et al. 2022.
\newblock Combining human and machine intelligence for clinical trial eligibility querying.
\newblock \textit{Journal of the American Medical Informatics Association} 29(7). 1161--1171.

\bibitem[Yuan et al. 2019]{yuan2019}
Yuan, Chi et al. 2019.
\newblock Criteria2Query: A natural language interface to clinical databases for cohort definition.
\newblock \textit{Journal of the American Medical Informatics Association} 26(4). 294--305.

\end{thebibliography}

\end{document}
